\documentclass[11pt,a4paper]{article}
% Préambule commun aux documents de PythonIPM/documentation.
% Compilation : tectonic <fichier>.tex
% tectonic compile avec XeTeX, qui lit l'UTF-8 nativement : surtout PAS inputenc ni fontenc,
% faits pour pdfLaTeX. Ce sont eux qui décodaient mal les caractères multi-octets — « cœur »
% ressortait en « cur », et les tirets cadratins en caractères parasites.
\usepackage{lmodern}
\usepackage[french]{babel}
\usepackage{microtype}
\usepackage[a4paper,margin=2.4cm]{geometry}
\usepackage{amsmath,amssymb}
\usepackage{booktabs}
\usepackage{longtable}
\usepackage{array}
\usepackage{ragged2e}
\usepackage{xcolor}
\usepackage{tcolorbox}
\usepackage{enumitem}
\usepackage{fancyhdr}
\usepackage[hidelinks]{hyperref}

\definecolor{encadre}{HTML}{F4F6F8}
\definecolor{bordure}{HTML}{4A6FA5}
\definecolor{alerte}{HTML}{B03A2E}

% Encadré « à retenir » : la lecture non technique du passage qui précède
\newtcolorbox{retenir}[1][]{
  colback=encadre, colframe=bordure, boxrule=0.6pt, arc=2pt,
  left=8pt, right=8pt, top=6pt, bottom=6pt,
  fonttitle=\bfseries, title={#1}
}

% Encadré d'avertissement : un point ouvert, un écart assumé
\newtcolorbox{attention}[1][]{
  colback=white, colframe=alerte, boxrule=0.6pt, arc=2pt,
  left=8pt, right=8pt, top=6pt, bottom=6pt,
  % le bandeau de titre a déjà la couleur d'alerte en fond : un titre `coltitle=alerte`
  % écrirait du rouge sur du rouge. On garde le blanc par défaut, comme l'encadré « retenir ».
  fonttitle=\bfseries, title={#1}
}

\newcommand{\col}[1]{\texttt{#1}}
% \ensuremath et non \textsubscript : \ipm est employé aussi bien dans le texte que dans les
% formules, et « M\textsubscript{0} » placé en mode mathématique est une erreur — silencieuse en
% PDF, mais qui fait échouer la conversion vers Word.
\newcommand{\ipm}{\ensuremath{M_{0}}}

\setlist[itemize]{itemsep=2pt, topsep=4pt}
\setlist[description]{style=nextline, leftmargin=1.2cm, itemsep=5pt}

% La source du document apparaît en tête de chaque page. Les documents RGPH la redéfinissent
% par \renewcommand{\source}{...} AVANT \begin{document}.
\newcommand{\source}{EHCVM 2021}

\pagestyle{fancy}
\fancyhf{}
\fancyhead[L]{\small\itshape IPM Côte d'Ivoire --- \source}
\fancyhead[R]{\small\thepage}
\renewcommand{\headrulewidth}{0.3pt}

\renewcommand{\arraystretch}{1.15}


\title{\bfseries Note sur la dimension Emploi\\[4pt]
\large Ce que mesure l'indicateur de chômage, et ce qu'il laisse de côté}
\author{Pipeline \col{PythonIPM} --- IPM Côte d'Ivoire, EHCVM 2021}
\date{Septembre 2026}

\begin{document}
\maketitle

\begin{abstract}
\noindent La dimension Emploi de l'IPM ivoirien repose sur deux indicateurs de poids égal, dont
l'un --- le chômage --- ne prive que 6,5~\% de la population et n'apporte que 1,6~\% de l'\ipm{}.
Cette note documente précisément ce que recouvre la formule « au chômage au sens du BIT » telle
qu'elle est implémentée, met au jour cinq décisions que l'énoncé ne mentionne pas, puis distingue
les trois notions que l'on confond couramment --- le chômage BIT, la sous-utilisation de la
main-d'\oe uvre (SU1 à SU4) et le NEET. Elle chiffre chacune sur l'EHCVM, mesure l'effet
de deux montages alternatifs désormais implémentés dans le pipeline (options \col{-{}-su1\_su3}
et \col{-{}-su3}), et conclut sur les arbitrages ouverts.
\end{abstract}

\section{Le point de départ}

Le tableau de référence national place dans la dimension Emploi deux indicateurs, à
\textbf{0,125} chacun :

\begin{description}
  \item[Chômage] « Un membre du ménage âgé de 17-40 ans est au chômage » --- au sens du BIT.
  \item[Emploi agricole de subsistance] « Le chef de ménage est occupé mais son activité se
  limite à l'agriculture sur son propre champ, sans salariat, apprentissage ni commerce ».
\end{description}

Leurs rendements sont sans commune mesure. À poids identique, l'emploi de subsistance prive
40,5~\% de la population et apporte \textbf{16,9~\%} de l'\ipm{} ; le chômage prive 6,5~\% et
apporte \textbf{1,6~\%}. Le second indicateur porte à lui seul 91~\% de la dimension. Cette note
explique pourquoi.

\section{Ce que « chômage au sens du BIT » recouvre exactement}

La construction se fait en trois temps : un test individuel à trois critères, un filtre d'âge,
puis une règle d'agrégation au ménage.

\subsection{Le test individuel}

Au niveau de chaque membre du ménage, dans
\col{01\_preconstruction\_matrice\_situationnelle.py} :

\begin{center}
\small
\begin{tabular}{@{}p{2.2cm}p{7.4cm}p{2.6cm}@{}}
\toprule
\textbf{Critère BIT} & \textbf{Questions EHCVM mobilisées} & \textbf{Condition} \\
\midrule
Sans emploi & 4.10 « a travaillé $\ge$ 1 h les 7 derniers jours » (qui résume 4.06 champ,
4.07 commerce, 4.08 salariat, 4.09 apprentissage) et 4.11 « possède un emploi mais en était
absent » & ni l'un ni l'autre \\
En recherche & 4.15 ou 4.17 « a cherché un emploi rémunéré au cours des 30 derniers jours » &
oui explicite \\
Disponible & 4.20 « quand sera-t-il disponible » $\in$ \{immédiatement ; dans 15 jours ;
entre 15 jours et un mois\} ou 4.19 « est néanmoins disponible tout de suite » & oui explicite \\
\bottomrule
\end{tabular}
\end{center}

Les trois conditions sont \textbf{cumulatives}. L'entonnoir, sur les 64\,491 individus de
l'enquête : \textbf{45\,808 sans emploi} (71,0~\%), dont \textbf{999 en recherche} (2,2~\%), dont
\textbf{740 disponibles} (74,1~\%).

\subsection{Le filtre d'âge et l'agrégation}

La tranche retenue est \textbf{17-40 ans}. L'âge est l'âge déclaré (question 1.04a, renseigné pour
6~\% des membres), à défaut l'année d'enquête moins l'année de naissance. Il reste
\textbf{561 chômeurs BIT sur 19\,814 personnes de 17-40 ans}, soit 2,8~\%.

Le ménage est déclaré privé si \textbf{au moins un} de ses membres de 17-40 ans est chômeur BIT :
\textbf{476 ménages sur 12\,965}, soit 3,7~\% des ménages et \textbf{6,5~\% de la population}
--- l'écart entre les deux tient à ce que les ménages concernés sont plus grands que la moyenne.

\begin{retenir}[La chaîne complète en une phrase]
Est privé le ménage qui compte au moins une personne de 17 à 40 ans qui n'a pas travaillé une
seule heure la semaine passée, \emph{et} a fait une démarche de recherche dans les 30 derniers
jours, \emph{et} se déclare disponible sous un mois.
\end{retenir}

\section{Cinq décisions que l'énoncé ne mentionne pas}

\begin{description}
  \item[Un ménage sur cinq ne peut structurellement pas être privé.] \textbf{2\,397 ménages}
  (18,5~\%) n'ont aucun membre de 17-40 ans. Ils sont classés \emph{non privés}, en application de
  la convention générale du pipeline (« un ménage non concerné n'est pas privé »). C'est un choix
  cohérent, mais pas neutre : un ménage de personnes âgées est hors d'atteinte de cet indicateur.
  À noter que le champ \col{eligibilite = membres\_17\_40}, déclaré dans la définition de
  l'indicateur, n'est \textbf{jamais utilisé dans le calcul} --- il ne figure que dans les
  journaux. Il documente une intention qui n'est pas implémentée.

  \item[Le traitement des valeurs manquantes est asymétrique.] Le critère « sans emploi » s'écrit
  en négation (\col{!= 1}), de sorte qu'une valeur manquante y est comptée comme un « oui » :
  \textbf{9\,575 individus} sont réputés sans emploi alors que 4.10 ou 4.11 n'est pas renseigné
  chez eux. Les critères de recherche et de disponibilité, eux, exigent un « oui » explicite. Les
  trois conditions étant cumulatives, la permissivité du premier est intégralement neutralisée par
  la sévérité des deux autres et le résultat final est juste --- mais l'équilibre tient à
  l'enchaînement, non à une décision explicite.

  \item[La disponibilité retenue est plus large que le standard BIT.] Le BIT demande une
  disponibilité sous \emph{deux semaines}. Le code accepte aussi la modalité 3 de la question 4.20,
  « entre 15 jours et un mois ».

  \item[Un chercheur sur quatre est écarté sur la disponibilité] (999 $\rightarrow$ 740). Pour qui
  a répondu oui à 4.15 ou 4.17, la disponibilité ne peut être établie que par 4.20, la question
  4.19 n'étant posée qu'aux non-chercheurs. Une non-réponse à 4.20 fait sortir la personne du
  chômage. C'est le point où la mesure est la plus restrictive.

  \item[La borne de 40 ans] exclut par construction tout chômeur de 41 ans et plus.
\end{description}

\section{Trois notions distinctes, trois dénominateurs}

On parle couramment de « trois définitions du chômage ». Il n'y en a qu'une. Autour d'elle
gravitent une \emph{famille} de mesures de sous-utilisation de la main-d'\oe uvre, et un
indicateur de jeunesse qui ne relève pas du marché du travail. Ce qui les sépare n'est pas le
numérateur, c'est le \textbf{dénominateur}.

\subsection{Le cadre : cinq blocs}

\begin{center}
\small
\begin{tabular}{@{}lrr@{}}
\toprule
\textbf{Bloc} & \textbf{Effectif} & \textbf{Part} \\
\midrule
Occupés & 18\,468 & 51,9~\% \\
Chômeurs BIT & 674 & 1,9~\% \\
\cmidrule(r){1-3}
\emph{\quad = main-d'\oe uvre} & \emph{19\,142} & \emph{53,7~\%} \\
\addlinespace
Cherche mais n'est pas disponible & 246 & 0,7~\% \\
Est disponible mais ne cherche pas & 572 & 1,6~\% \\
\cmidrule(r){1-3}
\emph{\quad = main-d'\oe uvre potentielle} & \emph{818} & \emph{2,3~\%} \\
\addlinespace
Autres inactifs & 15\,655 & 44,0~\% \\
\midrule
\textbf{Population de 15 ans et plus} & \textbf{35\,615} & \textbf{100~\%} \\
\bottomrule
\end{tabular}
\end{center}

\subsection{SU1 --- le taux de chômage}

\[ \mathrm{SU1} = \frac{\text{chômeurs}}{\text{occupés} + \text{chômeurs}}
   = \frac{674}{19\,142} = 3{,}5~\% \]

Il répond à la question : \emph{parmi ceux qui sont sur le marché du travail, combien n'y trouvent
pas de place ?} Deux mécanismes le tirent vers le bas dans une économie comme celle de la Côte
d'Ivoire, l'un au numérateur, l'autre au dénominateur.

Au \textbf{numérateur}, le critère de recherche active : renoncer, c'est sortir du chômage. Là où
chercher coûte du transport et où l'offre institutionnelle est peu accessible, l'abandon fait
disparaître statistiquement le problème.

Au \textbf{dénominateur}, et c'est le plus lourd, le critère d'\emph{une heure} : est occupé
quiconque a travaillé au moins une heure dans la semaine, y compris sur son propre champ et sans
rémunération. Sur l'EHCVM, \textbf{8\,135 personnes --- 44,0~\% des occupés --- n'ont fait que du
champ}, rien d'autre.

\begin{retenir}[Ce que SU1 dit et ne dit pas]
Un SU1 de 3,5~\% ne signifie pas que le marché du travail se porte bien. Il signifie que peu de
personnes se trouvent dans la configuration précise « sans aucune activité \emph{et} en recherche
\emph{et} disponible ». Dans une économie à dominante agricole et informelle, cette configuration
est rare par construction.
\end{retenir}

\subsection{SU3 --- chômage et main-d'\oe uvre potentielle}

Le BIT (19\textsuperscript{e} CIST, 2013) définit quatre mesures de sous-utilisation, SU1 à SU4
(LU1 à LU4 en anglais), en cercles concentriques : SU1 le chômage ; SU2 y ajoute le sous-emploi
lié au temps ; SU3 y ajoute la main-d'\oe uvre potentielle ; SU4 réunit les trois.

La \textbf{main-d'\oe uvre potentielle} rassemble ceux qui veulent travailler sans réunir les trois
critères, en deux branches symétriques : ceux qui \emph{cherchent sans être disponibles} (246) et
ceux qui \emph{sont disponibles sans chercher} (572) --- c'est là que loge le découragement.

\[ \mathrm{SU3} = \frac{\text{chômeurs} + \text{main-d'\oe uvre potentielle}}
                       {\text{main-d'\oe uvre} + \text{main-d'\oe uvre potentielle}}
   = \frac{674 + 818}{19\,142 + 818} = \frac{1\,492}{19\,960} = 7{,}5~\% \]

Le \textbf{dénominateur change lui aussi} : SU3 n'est pas SU1 avec un numérateur élargi, mais une
reconstruction des deux termes. On ne passe donc pas de l'un à l'autre par une règle de trois.

\begin{center}
\small
\begin{tabular}{@{}lrrr@{}}
\toprule
\textbf{Champ} & \textbf{SU1} & \textbf{SU3} & \textbf{Rapport} \\
\midrule
15 ans et plus & 3,5~\% & 7,5~\% & $\times$ 2,1 \\
17-40 ans (bande de l'IPM) & 5,5~\% & 11,3~\% & $\times$ 2,0 \\
\bottomrule
\end{tabular}
\end{center}

\subsection{Qui compose la main-d'\oe uvre potentielle}

La question 4.18 (raison principale de non-recherche) identifie les 572 personnes de la seconde
branche :

\begin{center}
\small
\begin{tabular}{@{}lrr@{}}
\toprule
\textbf{Raison de non-recherche (4.18)} & \textbf{Effectif} & \textbf{Part} \\
\midrule
Ménagère & 185 & 32,3~\% \\
\textbf{Manque d'emploi} & 102 & 17,8~\% \\
Attend le démarrage de sa propre entreprise & 81 & 14,2~\% \\
\textbf{Ne sait pas comment chercher} & 61 & 10,7~\% \\
Étudiant / Élève & 56 & 9,8~\% \\
Attend la réponse à une demande d'emploi & 50 & 8,7~\% \\
Ne veut pas travailler & 21 & 3,7~\% \\
Autres (maladie, âge, retraite, handicap) & 16 & 2,8~\% \\
\bottomrule
\end{tabular}
\end{center}

Les deux lignes en gras --- 163 personnes, 28~\% de la branche --- constituent le
\textbf{découragement au sens propre} : disponibles, mais convaincus qu'il n'y a rien à trouver ou
ne sachant pas s'y prendre. SU1 les classe parmi les inactifs, au même titre qu'un rentier.

Deux réserves sur cette liste. Les ménagères et les étudiants y figurent parce qu'ils ont déclaré
être disponibles --- ce qui est conforme au BIT, qui n'exclut personne sur le statut, mais pèse
42~\% de la branche. Et les 21 « ne veut pas travailler » sont en contradiction avec leur propre
déclaration de disponibilité : ils seraient à retirer lors de la construction.

\subsection{NEET --- un objet différent}

\emph{Not in Employment, Education or Training} : les 15-24 ans qui ne sont ni en emploi, ni en
études, ni en formation. C'est l'indicateur \textbf{ODD 8.6.1}. Son dénominateur n'est pas la
population active mais \textbf{l'ensemble des jeunes}. Un NEET peut être chômeur, découragé,
femme au foyer ou malade : l'indicateur ne mesure pas une offre de travail non satisfaite, il
mesure une \textbf{déconnexion} --- des jeunes qui n'accumulent de capital humain ni par l'école
ni par l'emploi.

\begin{center}
\small
\begin{tabular}{@{}lll r@{}}
\toprule
& \textbf{Question posée} & \textbf{Dénominateur} & \textbf{EHCVM} \\
\midrule
SU1 / BIT & Qui, sur le marché, n'y trouve pas de place ? & main-d'\oe uvre & 3,5~\% \\
SU3 & Qui veut travailler sans y parvenir ? & + main-d'\oe uvre potentielle & 7,5~\% \\
NEET & Quels jeunes ne construisent rien ? & population des 15-24 ans & $\approx$ 40~\% \\
\bottomrule
\end{tabular}
\end{center}

\section{Ce qui est constructible sur l'EHCVM}

\subsection{SU3 : oui, intégralement}

Toutes les composantes sont présentes. Les taux de réponse apparemment faibles de 4.15 (24~\%),
4.19 (52~\%) et 4.18 (57~\%) chez les 15-24 ans relèvent du \textbf{routage du questionnaire} et
non de la non-réponse : 4.19 n'est posée qu'aux non-chercheurs, par construction. Les deux
branches de la main-d'\oe uvre potentielle se reconstituent sans approximation, et 4.18 permet
d'isoler proprement les découragés.

\subsection{SU2 et SU4 : non}

Le sous-emploi lié au temps demande, outre les heures travaillées --- disponibles via 4.36
(jours par mois) et 4.37 (heures par jour) ---, la question « souhaitiez-vous travailler
davantage ? ». Elle ne figure pas dans le questionnaire. SU2 et SU4 sont donc hors de portée.

\subsection{NEET : partiellement}

\begin{description}
  \item[L'emploi] est renseigné à 100~\% par 4.10.
  \item[L'éducation] pose une difficulté de vague. La question 2.08a (scolarisé en 2021/22)
  n'existe \textbf{qu'en vague 2} --- 0~\% de réponses en vague 1. Il faut passer par 2.12
  (scolarisé en 2020/21), présente dans les deux vagues à 66~\%. Vérification faite, les valeurs
  manquantes de 2.12 correspondent \emph{exactement} aux 3\,320 jeunes n'ayant jamais fréquenté
  l'école (question 2.03), à 17 cas résiduels près : ils sont donc non scolarisés par déduction,
  et la couverture est en réalité complète. Le décalage d'une année scolaire, lui, demeure.
  \item[La formation] n'est pas mesurée. Les questions 2.05 et 2.06 identifient bien une
  « formation professionnelle » non formelle, mais au passé et sans période de référence : c'est
  un \emph{a déjà suivi}, non un \emph{suit actuellement}. Et l'apprentissage (4.09) est déjà
  compté comme un emploi.
\end{description}

\begin{attention}[Un NEET EHCVM serait en réalité un « NEE »]
Ni en emploi, ni en études --- sans le volet formation. Sur cette base approchée, 4\,079 jeunes
sur 10\,078 sont concernés, soit \textbf{40,5~\% des 15-24 ans}. Le chiffre surestime le NEET
d'autant que des jeunes sont en formation professionnelle en cours ; les questions 2.05 et 2.06
permettent d'en majorer le biais, à défaut de le corriger.
\end{attention}

\section{SU3 mis en \oe uvre : l'effet mesuré}

Les deux montages ont été calculés et sont disponibles dans le pipeline, par les options
\col{-{}-su1\_su3} (SU3 \emph{en plus} du chômage BIT, l'Emploi passant à trois indicateurs de
0,0833) et \col{-{}-su3} (SU3 \emph{à la place} du chômage BIT, à 0,125). Sans option, l'IPM
national publié est inchangé --- ce que le pipeline vérifie à chaque exécution.

Au niveau de l'indicateur, le passage de SU1 à SU3 fait passer le nombre de personnes concernées
de \textbf{561 à 1\,200} sur la bande 17-40 (2,8~\% $\rightarrow$ 6,0~\%), et le nombre de ménages
privés de \textbf{476 à 1\,027} (3,7~\% $\rightarrow$ 7,9~\% des ménages ; 6,5~\%
$\rightarrow$ \textbf{12,1~\%} de la population).

\begin{center}
\small
\begin{tabular}{@{}lrrrrr@{}}
\toprule
\textbf{Variante} & \textbf{$H$} & \textbf{$A$} & \textbf{\ipm{}} & \textbf{Vulnérables} &
\textbf{Sévère} \\
\midrule
IPM national (SU1) & 58,0~\% & 0,4940 & 0,2863 & 22,8~\% & 26,1~\% \\
SU1 + SU3 (\col{-{}-su1\_su3}) & 57,6~\% & 0,4750 & 0,2736 & 24,0~\% & 20,9~\% \\
SU3 seul (\col{-{}-su3}) & 59,3~\% & 0,4970 & 0,2949 & 22,7~\% & 27,1~\% \\
\bottomrule
\end{tabular}
\end{center}

\begin{center}
\small
\begin{tabular}{@{}lrrr@{}}
\toprule
\textbf{Contribution à l'\ipm{}} & \textbf{National} & \textbf{SU1 + SU3} & \textbf{SU3 seul} \\
\midrule
Chômage BIT (SU1) & 1,6~\% & 1,4~\% & --- \\
Sous-utilisation SU3 & --- & 2,4~\% & 3,2~\% \\
Emploi de subsistance & 16,9~\% & 11,2~\% & 16,5~\% \\
\midrule
\textbf{Dimension Emploi} & \textbf{18,5~\%} & \textbf{15,0~\%} & \textbf{19,7~\%} \\
\bottomrule
\end{tabular}
\end{center}

\subsection{Lecture}

\textbf{SU3 seul} est le montage cohérent. Il fait ce qu'on attendait : l'indicateur de chômage
double sa contribution (1,6~\% $\rightarrow$ 3,2~\%) sans bouleverser l'indice, qui monte de
0,2863 à \textbf{0,2949}. L'écart reste dans l'intervalle de confiance de l'IPM national
(\mbox{[0,2701 ; 0,3026]}) : le changement est cohérent avec la mesure actuelle, il ne la
contredit pas. Le rapport rural/urbain se resserre légèrement, de 2,85 à 2,72 --- SU3 capte du
découragement en ville, là où le chômage BIT strict est déjà le plus visible.

\textbf{SU1 + SU3} pose en revanche un problème de méthode, visible dans les chiffres. SU3
\emph{englobe} SU1 : tout ménage privé au titre du chômage BIT l'est aussi au titre de SU3, et se
voit donc compter deux fois la même privation. Le montage revient à donner un poids de
$2 \times 0{,}0833 = 0{,}167$ aux ménages à chômeur BIT et de 0,0833 aux autres --- une
pondération qu'aucun cadrage n'a décidée. Le symptôme est net : la dimension Emploi
\emph{perd} du poids (18,5~\% $\rightarrow$ 15,0~\%) parce que l'ajout d'un troisième indicateur
dilue l'emploi de subsistance, et la pauvreté sévère chute de 26,1~\% à 20,9~\% sans qu'aucune
réalité n'ait changé.

\begin{attention}[Si l'on veut deux indicateurs d'emploi non redondants]
Le montage propre n'est pas « SU1 + SU3 » mais \textbf{« SU1 + main-d'\oe uvre potentielle »} :
le second indicateur ne retiendrait alors que les personnes que SU1 ne voit pas --- celles qui
cherchent sans être disponibles et celles qui sont disponibles sans chercher. Les deux
indicateurs seraient disjoints et chaque ménage ne serait compté qu'une fois. Cette variante n'est
pas implémentée à ce jour ; elle s'obtiendrait en remplaçant \col{ind.su3} par
\col{ind.main\_oeuvre\_potentielle} à l'étape 01.
\end{attention}

\section{Les arbitrages ouverts}

\begin{enumerate}
  \item \textbf{Retenir SU3 seul plutôt que SU1.} C'est la piste la plus directe pour rendre à la
  dimension un équilibre entre ses deux jambes, et la seule des deux qui ne compte pas deux fois
  les mêmes ménages.

  \item \textbf{Réexaminer la borne de 40 ans et le seuil de disponibilité.} La première exclut
  les chômeurs plus âgés ; le second, en acceptant la modalité « entre 15 jours et un mois »,
  s'écarte déjà du standard BIT --- autant le trancher explicitement.

  \item \textbf{Implémenter le champ d'éligibilité, ou le retirer.} En l'état, il est déclaré mais
  inopérant. Le conserver sans l'appliquer entretient une ambiguïté sur le traitement des
  18,5~\% de ménages sans membre de 17-40 ans.

  \item \textbf{Ne pas substituer le NEET au chômage.} Il mesure autre chose --- un risque de
  pauvreté future plutôt qu'une privation présente --- et resterait incomplet faute du volet
  formation. S'il devait entrer dans l'IPM, ce serait comme indicateur supplémentaire assumé, non
  comme un chômage élargi.
\end{enumerate}

\begin{retenir}[SU3 et emploi de subsistance sont complémentaires, pas redondants]
SU3 corrige le \emph{numérateur} du chômage : il récupère ceux que le critère de recherche
écartait. Il ne touche pas au dénominateur --- les 8\,135 personnes qui ne font que du champ
restent comptées comme occupées. C'est précisément ce que l'indicateur d'emploi agricole de
subsistance capte déjà, en tenant lieu du SU2 que l'enquête ne permet pas de construire. Les deux
indicateurs de la dimension attaquent donc deux bouts opposés du même problème, et le
renforcement de l'un n'empiète pas sur l'autre.
\end{retenir}

Comme pour la dimension Santé, ces arbitrages relèvent du cadrage et non de la technique : le
critère de chômage se modifie en quelques lignes dans
\col{pipeline/01\_preconstruction\_matrice\_situationnelle.py}, et la règle du ménage dans
\col{pipeline/02\_construction\_matrice\_situationnelle.py}.

\vfill
\noindent\rule{\linewidth}{0.4pt}\\[2pt]
\small
Documents associés : \emph{Note sur la dimension Santé}, \emph{Méthodologie de calcul} et
\emph{Dictionnaire des sorties}.

\end{document}
